% !PRINTED 534
\weberpage{534}

%<<CONT
gleichungen gehören, die \gesperrt{Abel} allgemein auflösen gelehrt hat, und die wir also nach ihm \gesperrt{Abel'sche Gleichungen} nennen wollen\footnote{\gesperrt{Abel}, Mémoire sur une classe d'équations résolubles algébriquement. Crelle's Journal f. Mathematik, Bd.~4, 1829. Oeuvres complètes, nouvelle édition, 1881, Bd.~1, S.~418.}.

\gesperrt{Eine Gleichung $m^{\mathrm{ten}}$ Grades $F(x) = 0$ mit den Wurzeln $\alpha, \alpha_1, \alpha_2 \ldots \alpha_{m-1}$ heisst eine Abel'sche Gleichung, wenn jede Wurzel rational durch eine von ihnen, $\alpha$, ausdrückbar ist, und wenn, falls}
\weq{2}{\alpha_1 = \Theta_1(\alpha), \quad \alpha_2 = \Theta_2(\alpha) \ldots \alpha_{m-1} = \Theta_{m-1}(\alpha)}
\gesperrt{diese rationalen Ausdrücke sind, die Bedingung}
\weq{3}{\Theta_h\Theta_k(\alpha) = \Theta_k\Theta_h(\alpha)}
\gesperrt{für je zwei dieser Functionen besteht.}

Es bedeutet hierin das Zeichen $\Theta_h\Theta_k(\alpha)$, dass die Function $\Theta_h(x)$ für das Argument $x = \Theta_k(\alpha)$ gebildet werden soll. Selbstverständlich bezieht sich diese ganze Definition auf einen bestimmten als rational angenommenen Körper $\Omega$. Zu diesen Gleichungen gehören gewiss, wenn $\Omega$ der Körper der rationalen Zahlen ist, die Gleichungen, durch die die Einheitswurzeln bestimmt sind. Denn jede $m^{\mathrm{te}}$ Einheitswurzel ist, wenn $r$ eine primitive unter ihnen ist, in der Form $\Theta_h(r) = r^h$ enthalten, und es ist
\[
\Theta_h\Theta_k(r) = \Theta_k\Theta_h(r) = r^{hk}.
\]
Ebenso gehört, wenn der Körper $\Omega$ die $m^{\mathrm{ten}}$ Einheitswurzeln $r$ enthält, die \gesperrt{reine Gleichung}
\[
x^m - a = 0,
\]
in der $a$ dem Körper $\Omega$ angehört, zu den \gesperrt{Abel'schen Gleichungen}. Denn ist $\alpha$ eine ihrer Wurzeln, so sind alle Wurzeln in der Form
\[
\alpha_h = r^h\alpha
\]
enthalten; es ist also $\Theta_h(x) = r^hx$ und daher
\[
\Theta_h\Theta_k(x) = \Theta_k\Theta_h(x) = r^{h+k}x.
\]
Wenn die Function $F(x)$ nicht irreducibel ist, so hat sie einen bestimmten irreducibeln Factor $\varphi(x)$, der die Wurzel $\alpha$ hat, und unter den Wurzeln von $\varphi(x) = 0$ bestehen gleichfalls die durch (2) und (3) ausgesprochenen Relationen. Es ist daher $\varphi(x) = 0$ nach \S.\,145 eine \gesperrt{Galois'sche} Resolvente von
%CONT>>
